% 04-components.tex
\chapter{Subsystems and Components}
\label{ch:components}

\section{Image Management}
\label{sec:image-management}

Doki implements a complete OCI-compliant image management system. Images
are stored in a layout compatible with the Docker image specification,
using content-addressable storage for layers and manifests.

\subsection{Image Storage Layout}
\label{sec:image-layout}

Images are stored under \texttt{\$DOKI\_DATA\_DIR/images/} with the
following structure:

\begin{codebox}[title={Image storage directory layout}]{text}
images/
  manifests/
    <sha256>.json          -- Image manifest
  layers/
    <sha256>/              -- Extracted layer contents
  blobs/
    sha256/
      <sha256>             -- Compressed layer tarballs
\end{codebox}

Each manifest contains the image configuration (entrypoint, cmd, env,
labels) and a list of layer digests. Layers are extracted sequentially
into the \texttt{layers/} directory, with later layers overwriting
earlier ones to produce the final filesystem.

\subsection{Image Pull Process}
\label{sec:image-pull}

The image pull process follows the OCI distribution specification:

\begin{figure}[htbp]
  \centering
  \begin{tikzpicture}[
    node distance=0.4cm,
    every node/.style={font=\sffamily\footnotesize},
    stage/.style={
      rectangle, rounded corners=2pt,
      draw=nodeblue!60, fill=fillblue,
      minimum width=2.0cm, minimum height=0.8cm,
      text=textdark, align=center, inner sep=4pt
    }
  ]
    \node[stage] (resolve) {Resolve\\Tag};
    \node[stage, right=0.8cm of resolve] (manifest) {Fetch\\Manifest};
    \node[stage, right=0.8cm of manifest] (layers) {Download\\Layers};
    \node[stage, right=0.8cm of layers] (extract) {Extract\\Layers};
    \node[stage, right=0.8cm of extract] (config) {Apply\\Config};

    \draw[arr] (resolve) -- (manifest);
    \draw[arr] (manifest) -- (layers);
    \draw[arr] (layers) -- (extract);
    \draw[arr] (extract) -- (config);

    % Parallel layer downloads
    \draw[arrDash] ([yshift=-0.5cm]layers.south) -- ++(0,-0.4)
      node[below, font=\sffamily\tiny, text=textgray] {Parallel download};
  \end{tikzpicture}
  \caption{Image pull pipeline. Layers are downloaded in parallel
    when multiple layers are present.}
  \label{fig:image-pull}
\end{figure}

\subsection{Search and Registry Interaction}
\label{sec:search}

Doki can search Docker Hub and private registries. The search command
parses the Docker Hub v2 API response format, which uses
\texttt{repo\_name} and \texttt{short\_description} as JSON keys.
The \texttt{History} struct uses PascalCase JSON tags
(\texttt{Created} as \texttt{int64}, \texttt{CreatedBy},
\texttt{Size}, \texttt{Comment}) to match the daemon response format.

The search flow consists of three stages:

\begin{enumerate}[leftmargin=*, itemsep=3pt]
  \item \textbf{API request.} The CLI sends a GET request to the
        registry's search endpoint with the query term as a parameter.
  \item \textbf{Response parsing.} The JSON response is decoded into
        a list of repository structs containing name, description,
        star count, and pull count.
  \item \textbf{Display formatting.} Results are formatted as a table
        with truncated descriptions and sorted by relevance.
\end{enumerate}

\section{Container Runtime}
\label{sec:container-runtime}

The container runtime manages process execution within isolated
environments. It supports multiple isolation backends through a
pluggable runner architecture.

\subsection{Proot Runner}
\label{sec:proot-runner}

On Android/Termux, the proot runner is the primary isolation mechanism.
It uses \texttt{ptrace()} to intercept system calls and provides:

\begin{itemize}[leftmargin=*, itemsep=3pt]
  \item Virtual filesystem views (bind mounts without root)
  \item Process ID namespace virtualization
  \item User ID mapping
  \item Hostname and DNS configuration
\end{itemize}

The proot binary is located using \texttt{FindProotBinary()}, which
searches \texttt{\$PATH} and falls back to Doki's bundled
\texttt{doki-proot}. Environment variables \texttt{LD\_PRELOAD} and
\texttt{LD\_LIBRARY\_PATH} are automatically stripped from the proot
environment to prevent library injection attacks.

\subsection{Namespace Runner}
\label{sec:namespace-runner}

On systems with root access, the namespace runner provides stronger
isolation using kernel-level namespaces:

\begin{itemize}[leftmargin=*, itemsep=3pt]
  \item PID namespace (process isolation)
  \item Mount namespace (filesystem isolation)
  \item Network namespace (network isolation)
  \item UTS namespace (hostname isolation)
  \item IPC namespace (inter-process communication isolation)
  \item User namespace (user ID mapping without root)
\end{itemize}

\subsection{Native Runner}
\label{sec:native-runner}

The native runner executes processes with minimal isolation, using only
cgroups for resource limiting. It is the fastest runtime but provides
the weakest isolation. It is typically used in trusted environments
where performance is critical.

\begin{table}[htbp]
  \centering
  \caption{Runtime comparison: proot vs.\ native runner.}
  \label{tab:runtime-comparison}
  \begin{tabularx}{\textwidth}{l X X}
    \toprule
    \textbf{Feature} & \textbf{Proot Runner} & \textbf{Native Runner} \\
    \midrule
    Isolation level  & Medium (ptrace) & Low (cgroups only) \\
    Startup time     & 100--150\,ms    & $<$10\,ms \\
    Memory overhead  & 15\,MB          & $<$2\,MB \\
    Root required    & No              & Optional \\
    Use case         & Android, untrusted & Trusted, performance-critical \\
    \bottomrule
  \end{tabularx}
\end{table}

\section{Dokifile Builder}
\label{sec:builder}

Doki includes a complete Dockerfile-compatible builder that parses
Dockerfile instructions and builds OCI-compliant images.

\subsection{Supported Instructions}
\label{sec:supported-instructions}

The builder supports the following Dockerfile instructions:

\begin{table}[htbp]
  \centering
  \caption{Supported Dockerfile instructions.}
  \label{tab:dockerfile-instructions}
  \begin{tabularx}{\textwidth}{l X}
    \toprule
    \textbf{Instruction} & \textbf{Description} \\
    \midrule
    \texttt{FROM}         & Base image specification (supports multi-stage) \\
    \texttt{RUN}          & Execute shell commands during build \\
    \texttt{COPY}         & Copy files from build context to image \\
    \texttt{ADD}          & Copy files with auto-extraction for tarballs \\
    \texttt{WORKDIR}      & Set working directory \\
    \texttt{ENV}          & Set environment variables \\
    \texttt{EXPOSE}       & Document listening ports \\
    \texttt{CMD}          & Default command (overridable) \\
    \texttt{ENTRYPOINT}   & Fixed entrypoint command \\
    \texttt{USER}         & Set process user \\
    \texttt{LABEL}        & Add metadata labels \\
    \texttt{VOLUME}       & Declare mount points \\
    \texttt{ARG}          & Build-time arguments \\
    \texttt{ONBUILD}      & Trigger instructions for downstream images \\
    \bottomrule
  \end{tabularx}
\end{table}

\subsection{Layer Cache}
\label{sec:layer-cache}

The builder implements a layer cache that stores intermediate build
results. Each instruction produces a layer that is cached by its
content hash. On subsequent builds, unchanged layers are reused from
cache, significantly reducing build times.

The \texttt{extractBaseImageToDir} function handles the extraction of
base image layers, including gzip-compressed tarballs with symlinks.
Layer ordering is preserved to maintain filesystem semantics where
later layers override earlier ones.

\section{Compose Engine}
\label{sec:compose-engine}

The \texttt{doki-compose} binary provides Docker Compose compatibility,
allowing users to define multi-container applications using YAML
configuration files.

\subsection{Service Dependencies}
\label{sec:service-dependencies}

The compose engine resolves service dependencies using a directed
acyclic graph (DAG) traversal. Services are started in topological
order, with dependency conditions including:

\begin{itemize}[leftmargin=*, itemsep=3pt]
  \item \texttt{service\_started} --- wait for the dependency container
        to start
  \item \texttt{service\_healthy} --- wait for a health check to pass
        (parsed but not yet executed by the runtime)
  \item \texttt{service\_completed\_successfully} --- wait for the
        dependency container to exit with code 0
\end{itemize}

If a circular dependency is detected, the \texttt{orderServices}
function returns an error rather than silently randomizing the start
order.

\subsection{Profile Filtering}
\label{sec:profile-filtering}

Services with a \texttt{profiles} key are only included when one of
their profiles is activated via the \texttt{--profile} flag. This
allows defining optional services (such as debug tools or monitoring
stacks) that are not started by default.

\begin{figure}[htbp]
  \centering
  \begin{tikzpicture}[
    node distance=1.0cm,
    every node/.style={font=\sffamily\footnotesize},
    svc/.style={
      rectangle, rounded corners=2pt,
      draw=nodeblue!60, fill=fillblue,
      minimum width=2.0cm, minimum height=0.7cm,
      text=textdark, align=center, inner sep=4pt
    },
    arr/.style={
      -{Stealth[length=2.5mm, width=2mm]},
      draw=linegray, line width=0.5pt
    }
  ]
    \node[svc] (web) {web};
    \node[svc, right=1.2cm of web] (api) {api};
    \node[svc, right=1.2cm of api] (db) {db};
    \node[svc, below=0.8cm of api] (redis) {redis};
    \node[svc, below=0.8cm of web] (nginx) {nginx};

    \draw[arr] (web) -- (api);
    \draw[arr] (api) -- (db);
    \draw[arr] (api) -- (redis);
    \draw[arr] (nginx) -- (web);

    \node[annot, below=0.3cm of db] {arrows = depends-on};
  \end{tikzpicture}
  \caption{Example Compose service dependency graph. Services are
    started in topological order.}
  \label{fig:compose-dag}
\end{figure}
